% Generated draft: requires mathematical and source review.
\documentclass[11pt]{article}
\usepackage[margin=25mm]{geometry}
\usepackage{fontspec}
\setmainfont{Latin Modern Roman}
\usepackage{amsmath,amssymb,mathtools,unicode-math}
\setmathfont{Latin Modern Math}
\usepackage{enumitem,needspace,xurl,hyperref}
\setlength{\parindent}{0pt}
\setlength{\parskip}{5pt}
\setlength{\emergencystretch}{3em}
\newcommand{\N}{\mathbb{N}}
\newcommand{\Z}{\mathbb{Z}}
\newcommand{\Q}{\mathbb{Q}}
\newcommand{\R}{\mathbb{R}}
\newcommand{\C}{\mathbb{C}}
\newcommand{\F}{\mathbb{F}}
\newcommand{\A}{\mathbb{A}}
\newcommand{\PP}{\mathbb{P}}
\newcommand{\E}{\mathbb{E}}
\newcommand{\eps}{\varepsilon}
\newcommand{\dd}{\,\mathrm{d}}
\begin{document}
\section*{UnsolvedMath ID: 30006264}
\textbf{Source problem code:} catalogue entry
\subsection{Definitions and mathematical statement}
The supplied record concerns Bayesian inverse problems whose prior is a random-tree Besov prior. It does not specify the random-tree distribution, the associated Besov parameters, the parameter space, the observation model, or the class of inverse problems. Consequently, the following is a schematic formulation; these missing choices must be supplied before it becomes a fully quantified mathematical problem.

For an inverse-problem instance \(I\), let \((U_I,\mathcal U_I)\) be its measurable parameter space, let \(\mu_{0,I}\) be its random-tree Besov prior, and let \(L_I(y\mid u)\) be a nonnegative measurable likelihood for the observed data \(y\). In a likelihood-based formulation, provided
\[
0<Z_I(y):=\int_{U_I}L_I(y\mid u)\,\mu_{0,I}(\mathrm du)<\infty,
\]
the posterior is the probability measure
\[
\mu_I^y(A)=\frac{1}{Z_I(y)}\int_A L_I(y\mid u)\,\mu_{0,I}(\mathrm du),\qquad A\in\mathcal U_I.
\]
This displays the usual posterior construction and its normalization condition; the supplied record does not itself specify a likelihood or assumptions ensuring this condition.

\textbf{Sampling problem.} Develop general algorithms for sampling the posteriors arising from random-tree Besov priors. The supplied expanded statement additionally requests provable convergence and quantitative mixing and computational-complexity guarantees. To make that expanded target precise, one must specify an admissible class \(\mathcal I\) of instances, the allowed observations, a discrepancy \(D_I\) between probability measures, and a computational cost model. For such specified choices, a schematic convergence requirement for algorithmic output laws \(\nu_{I,y,n}\) is
\[
D_I(\nu_{I,y,n},\mu_I^y)\le B_{I,y}(n),\qquad
\lim_{n\to\infty}B_{I,y}(n)=0,
\]
for every admissible instance and observation, with an explicit bound and stated dependence on all relevant parameters and initialization. A corresponding complexity guarantee would bound the cost of producing an output law within discrepancy \(\varepsilon\) of \(\mu_I^y\), for each \(\varepsilon>0\), including any discretization, truncation, and sampling errors introduced by the algorithm.

For a Markov chain method, a mixing formulation would additionally specify a transition kernel \(P_{I,y}\) having invariant measure \(\mu_I^y\), admissible initial laws \(\lambda\), and explicit bounds of the form
\[
D_I(\lambda P_{I,y}^{\,n},\mu_I^y)\le B_{I,y,\lambda}(n)\longrightarrow0.
\]
This conditional formulation does not restrict all sought algorithms to Markov chain methods. No discrepancy, rate, uniformity over instances, initialization class, or complexity exponent is fixed by the supplied source. The original extracted question asks how to sample these posteriors; the quantitative guarantees appear in the expanded statement and are not explicit in that original question.
\subsection{Short English statement}
Develop sampling algorithms for Bayesian inverse-problem posteriors with random-tree Besov priors. The expanded registry statement asks for convergence, mixing, and complexity guarantees, but the prior family, inverse-problem assumptions, error criterion, and cost model remain unspecified. The supplied dated assessment labels the problem partially solved; this draft does not independently verify that assessment.
\subsection{Sources}
\begin{itemize}
\item Supplied local UnsolvedMath registry record, permanent ID \(30006264\), title: Posterior Sampling for Random-Tree Besov Priors; exact source problem code: \texttt{catalogue entry}. The record attributes the question to the Oberwolfach Reports workshop collection, gives proposed year 2025, and names no individual proposer.
\item Supplied source citation: \emph{Uncertainty Quantification} (2025), Oberwolfach Reports, DOI 10.4171/owr/2025/21, pp. 1011--1078. \url{https://doi.org/10.4171/owr/2025/21}. The supplied extracted question asks how to sample posteriors arising from random-tree Besov priors; no pinpoint page is supplied.
\item Supplied literature source, OWR 2025 report: \url{https://ems.press/content/serial-article-files/51850}.
\item Supplied literature source, Multilevel MCMC with random-tree Besov priors: \url{https://arxiv.org/abs/2302.00678}.
\item Supplied literature source, Random-tree Besov priors: \url{https://www.aimsciences.org/article/doi/10.3934/ipi.2022059}.
\item Literature assessment dated exactly \texttt{2026-08-22}: the supplied curation labels the problem partially solved and states that independence Metropolis and multilevel MCMC methods provide effective posterior sampling for important elliptic inverse problems with random-tree Besov priors, while a general sampling solution is not established. This is supplied dated curation derived from the OWR material, not an independently verified literature assessment; its time-sensitive status and cited claims require independent verification.
\end{itemize}
\end{document}
